% =====================================================================
%  Banco complementar --- Teorema da Superposicao  (REVISADO)
%  Substitui a versao gerada por template (9 clones em N2, 10 em N3 --
%  todos com o mesmo enunciado sobre potencia -- e 8 questoes
%  conceituais indevidamente marcadas como N4).
%  Sem sobreposicao com o cap06, que ja traz: como desativar fontes
%  ideais (MC), o que pode ser superposto (MC), V e I no resistor de
%  3 ohm com duas fontes, potencial com fonte de tensao + corrente,
%  o erro de somar potencias (analise do erro do estudante) e o
%  circuito com tres fontes independentes.
%  Distribuicao mantida: 6 N1 + 9 N2 + 10 N3 + 8 N4 = 33
% =====================================================================

\subsubsection*{Banco complementar --- Teorema da Superposição}

\begin{atencao}[Organização do banco]
A superposição decorre de uma única propriedade: em rede linear,
$\mathbf{v}=\mathbf{G}^{-1}\mathbf{b}$ é \emph{linear} no vetor de excitações.
Daí três consequências que organizam todo o tópico: tensões e correntes se
somam; \textbf{potências não}; e fontes \emph{dependentes} nunca são desativadas,
porque pertencem à rede, não à excitação. O N3 trata de fontes reais,
coeficientes de transferência e falha em circuitos não lineares; o N4 exige
demonstração, algoritmo e projeto de ensaio.
\end{atencao}

\blocofixacao

\begin{questao}[1]{}
Um nó $A$ é alimentado por uma fonte de $\SI{12}{\volt}$ através de
$\SI{6}{\ohm}$ e ligado ao terra por $\SI{3}{\ohm}$; há ainda uma fonte de
corrente ligada a $A$. Determine a contribuição \emph{apenas} da fonte de
tensão para $V_A$.
\resposta{$V_A'=\SI{4}{\volt}$}
\resolucao{Desativar a fonte de corrente significa substituí-la por um
\textbf{circuito aberto} --- o ramo simplesmente desaparece. Resta um divisor:
\[
V_A'=12\cdot\frac{3}{6+3}=\SI{4}{\volt}.
\]
Esta é apenas uma \emph{parcela}: o valor real de $V_A$ só sai depois de somada
a contribuição da outra fonte.}
\end{questao}

\begin{questao}[1]{}
No mesmo circuito, a fonte de corrente vale $\SI{2}{\ampere}$ e injeta em $A$.
Determine a contribuição \emph{apenas} dela.
\resposta{$V_A''=\SI{4}{\volt}$}
\resolucao{Desativar a fonte de tensão ideal significa substituí-la por um
\textbf{curto}. Os dois resistores ficam então em paralelo:
\[
R=\frac{6\cdot3}{9}=\SI{2}{\ohm}
\;\Longrightarrow\;
V_A''=2\cdot2=\SI{4}{\volt}.
\]
\textbf{Memorize a assimetria:} fonte de \emph{tensão} ideal $\to$ curto (pois
impõe $\SI{0}{\volt}$); fonte de \emph{corrente} ideal $\to$ aberto (pois impõe
$\SI{0}{\ampere}$).}
\end{questao}

\begin{questao}[1]{}
Some as duas contribuições anteriores e determine $V_A$.
\resposta{$V_A=\SI{8}{\volt}$}
\resolucao{
\[
V_A=V_A'+V_A''=4+4=\SI{8}{\volt}.
\]
\textbf{Verificação por análise nodal:}
\[
\frac{V_A-12}{6}+\frac{V_A}{3}=2
\;\Longrightarrow\;
V_A-12+2V_A=12
\;\Longrightarrow\;
V_A=\SI{8}{\volt}.\ \checkmark
\]
A concordância não é coincidência --- as duas técnicas são consequências da
mesma linearidade.}
\end{questao}

\begin{questao}[1]{}
Um circuito contém quatro fontes independentes. Quantas análises parciais a
superposição exige, e o que se faz em cada uma?
\resposta{Quatro análises; em cada uma, apenas uma fonte permanece ativa}
\resolucao{Uma análise por fonte independente: mantém-se \emph{uma} ativa e
desativam-se as outras três (tensão $\to$ curto, corrente $\to$ aberto). Ao
final, somam-se as quatro contribuições \emph{com sinal}.\\
\textbf{Observação de custo:} o esforço cresce \emph{linearmente} com o número
de fontes, mas cada análise parcial costuma ser bem mais simples que o circuito
completo. A vantagem se inverte quando há muitas fontes --- com quatro ou mais,
a análise nodal em geral é mais rápida.}
\end{questao}

\begin{questao}[1]{}
Duas fontes produzem, sobre o mesmo resistor, contribuições de
$\SI{+7}{\volt}$ e $\SI{-3}{\volt}$. Determine a tensão resultante e comente.
\resposta{$U=\SI{4}{\volt}$}
\resolucao{
\[
U=(+7)+(-3)=\SI{4}{\volt}.
\]
O sinal negativo indica que a segunda fonte produz, sozinha, uma tensão de
polaridade \emph{oposta} à adotada como referência. \textbf{Somar módulos seria
erro grave} --- daria $\SI{10}{\volt}$.\\
Fixe uma polaridade de referência \emph{antes} de começar e mantenha-a em todas
as análises parciais: é a única forma de os sinais saírem corretos
automaticamente.}
\end{questao}

\begin{questao}[1]{}
O teorema da superposição pode ser aplicado a um circuito que contenha um
resistor cuja resistência varia com a corrente?
\begin{alternativas}[2]
\item sim, sempre
\item sim, desde que haja apenas uma fonte
\item não, pois o circuito deixa de ser linear
\item não, pois a potência não se superpõe
\end{alternativas}
\resposta{(c)}
\resolucao{A superposição exige \textbf{linearidade}: a resposta deve ser
proporcional à excitação. Um resistor cuja resistência depende da corrente
quebra essa proporcionalidade, e a soma das respostas parciais deixa de valer.\\
A alternativa (d) enuncia um fato verdadeiro, mas irrelevante aqui --- a
potência não se superpõe nem mesmo em circuitos perfeitamente lineares, por uma
razão diferente (ela é \emph{quadrática}, não não linear).\\
E (b) é curiosa: com uma única fonte não há o que superpor, então a questão nem
se coloca.}
\end{questao}

\blocoaplicacao

\begin{questao}[2]{}
Um nó $A$ recebe $\SI{24}{\volt}$ através de $\SI{12}{\ohm}$,
$\SI{12}{\volt}$ através de $\SI{6}{\ohm}$, uma fonte de corrente de
$\SI{1}{\ampere}$, e liga-se ao terra por $\SI{4}{\ohm}$. Determine $V_A$ por
superposição.
\resposta{$V_A=\SI{10}{\volt}$}
\resolucao{A condutância total do nó é a mesma nas três análises:
\[
G=\frac1{12}+\frac16+\frac14=\SI{0.5}{\siemens}.
\]
\textbf{Só a fonte de $\SI{24}{\volt}$} (as outras: curto e aberto):
$V'=\dfrac{24/12}{0{,}5}=\SI{4}{\volt}$.\\
\textbf{Só a de $\SI{12}{\volt}$:} $V''=\dfrac{12/6}{0{,}5}=\SI{4}{\volt}$.\\
\textbf{Só a de corrente:} $V'''=\dfrac{1}{0{,}5}=\SI{2}{\volt}$.\\
\[
V_A=4+4+2=\SI{10}{\volt}.
\]
\textbf{Atalho revelado:} como todas as fontes de tensão são aterradas, $G$ não
muda entre as análises --- e a soma das três parcelas reproduz exatamente a
equação nodal única. Superposição e análise nodal coincidem passo a passo neste
caso.}
\end{questao}

\begin{questao}[2]{}
No circuito anterior, determine a contribuição \emph{percentual} de cada fonte
para $V_A$.
\resposta{$40\%$, $40\%$ e $20\%$}
\resolucao{
\[
\frac{4}{10}=40\%;\qquad \frac{4}{10}=40\%;\qquad \frac{2}{10}=20\%.
\]
\textbf{Utilidade prática:} a decomposição identifica qual fonte domina o
resultado. Se a fonte de corrente for a menos estável, seu erro afeta $V_A$
com peso de apenas $20\%$; se o objetivo for estabilizar $V_A$, o esforço deve
ir para as duas fontes de tensão.\\
Essa leitura --- impossível pela análise nodal direta, que entrega só o total ---
é a principal razão pela qual a superposição continua sendo usada mesmo quando
não é o método mais rápido.}
\end{questao}

\begin{questao}[2]{}
Duas fontes de corrente injetam $\SI{3}{\ampere}$ e retiram
$\SI{1}{\ampere}$ de um mesmo nó, ligado ao terra por
$\SI{10}{\ohm}\parallel\SI{15}{\ohm}$. Determine a tensão por superposição e
verifique.
\resposta{$V=\SI{12}{\volt}$}
\resolucao{A resistência do nó é $10\parallel15=\SI{6}{\ohm}$ em ambas as
análises (desativar uma fonte de corrente apenas abre seu ramo).
\[
V'=3(6)=\SI{18}{\volt};
\qquad
V''=-1(6)=\SI{-6}{\volt};
\qquad
V=18-6=\SI{12}{\volt}.
\]
\textbf{Verificação direta:} corrente líquida $3-1=\SI{2}{\ampere}$, logo
$V=2(6)=\SI{12}{\volt}$ \checkmark.\\
Com fontes de corrente no mesmo nó, a superposição é trivial --- elas se somam
algebricamente antes de qualquer conta. O método só ganha relevância quando as
fontes atuam em \emph{pontos diferentes} da rede.}
\end{questao}

\begin{questao}[2]{}
Duas fontes de tensão de $\SI{18}{\volt}$ e $\SI{6}{\volt}$, com polaridades
\textbf{opostas}, estão em série com $R=\SI{4}{\ohm}$. Determine a corrente por
superposição.
\resposta{$I=\SI{3}{\ampere}$}
\resolucao{Com a malha série, cada fonte age sozinha sobre o mesmo $R$ (a outra,
curto-circuitada, não altera a resistência):
\[
I'=\frac{18}{4}=\SI{4.5}{\ampere};
\qquad
I''=-\frac{6}{4}=\SI{-1.5}{\ampere};
\qquad
I=4{,}5-1{,}5=\SI{3}{\ampere}.
\]
\textbf{Verificação:} tensão líquida $18-6=\SI{12}{\volt}$, logo
$I=12/4=\SI{3}{\ampere}$ \checkmark.\\
O sinal negativo de $I''$ é essencial: a fonte de $\SI{6}{\volt}$ se opõe à de
$\SI{18}{\volt}$. Somar módulos daria $\SI{6}{\ampere}$ --- o dobro do
correto.}
\end{questao}

\begin{questao}[2]{}
Determine, por superposição, a \textbf{corrente} no resistor de
$\SI{4}{\ohm}$ do circuito de três fontes ($V_A=\SI{10}{\volt}$).
\resposta{$I=\SI{2.5}{\ampere}$}
\resolucao{As contribuições de tensão no nó já foram obtidas: $4$, $4$ e
$\SI{2}{\volt}$. Como o resistor liga $A$ ao terra,
\[
I'=\frac{4}{4}=1;\quad I''=\frac{4}{4}=1;\quad I'''=\frac{2}{4}=0{,}5
\;\Longrightarrow\;
I=\SI{2.5}{\ampere}.
\]
\textbf{Verificação:} $I=V_A/4=10/4=\SI{2.5}{\ampere}$ \checkmark.\\
\textbf{Observação:} correntes se superpõem tão livremente quanto tensões ---
ambas são \emph{lineares} nas excitações. É a mesma propriedade, aplicada a
outra grandeza.}
\end{questao}

\begin{questao}[2]{}
Em um circuito, a contribuição de $V_1$ para certa corrente é
$\SI{200}{\milli\ampere}$. Se $V_1$ for \textbf{triplicada}, mantidas as demais
fontes, o que ocorre com essa contribuição e com a corrente total?
\resposta{A contribuição triplica ($\SI{600}{\milli\ampere}$); a corrente total
aumenta em $\SI{400}{\milli\ampere}$}
\resolucao{Cada contribuição é \emph{proporcional} à sua fonte, portanto
$200\to\SI{600}{\milli\ampere}$.\\
As demais contribuições \textbf{não mudam}, pois nada foi alterado nas outras
fontes nem na rede. Logo a corrente total aumenta exatamente pelo acréscimo:
$600-200=\SI{400}{\milli\ampere}$.\\
\textbf{Poder do método:} para responder "o que acontece se eu mexer nesta
fonte", não é preciso resolver o circuito de novo --- basta escalar a parcela
correspondente. Essa é a base da análise de sensibilidade em redes lineares.}
\end{questao}

\begin{questao}[2]{}
Um circuito com duas fontes tem contribuições de $\SI{6}{\volt}$ e
$\SI{2}{\volt}$ sobre um resistor de $\SI{4}{\ohm}$. Calcule a potência de
duas formas: somando potências parciais e usando a tensão total. Compare.
\resposta{$\SI{10}{\watt}$ (errado) contra $\SI{16}{\watt}$ (correto)}
\resolucao{\textbf{Somando potências parciais (incorreto):}
\[
P'=\frac{6^{2}}{4}=9;\qquad P''=\frac{2^{2}}{4}=1;\qquad
P'+P''=\SI{10}{\watt}.
\]
\textbf{Pela tensão total (correto):}
\[
U=6+2=\SI{8}{\volt}
\;\Longrightarrow\;
P=\frac{8^{2}}{4}=\SI{16}{\watt}.
\]
A diferença é de $\SI{6}{\watt}$ --- erro de $37{,}5\%$ no resultado correto.\\
\textbf{Procedimento correto:} superponha \emph{tensões} (ou correntes), some, e
só então calcule a potência. Nunca superponha potências. A origem algébrica
dessa falha é explorada no bloco de desafio.}
\end{questao}

\begin{questao}[2]{}
Uma fonte \textbf{real} de $\SI{12}{\volt}$ com $r=\SI{2}{\ohm}$ e outra de
$\SI{6}{\volt}$ com $r=\SI{3}{\ohm}$ alimentam em paralelo uma carga de
$\SI{6}{\ohm}$. Determine a tensão na carga por superposição.
\resposta{$U=\SI{8}{\volt}$}
\resolucao{\textbf{Ponto crítico:} desativar uma fonte real significa zerar sua
\textbf{f.e.m.}, mantendo sua resistência interna no circuito.\\
Condutância total (a mesma nas duas análises):
$G=\frac12+\frac13+\frac16=\SI{1}{\siemens}$.
\[
U'=\frac{12/2}{1}=\SI{6}{\volt};
\qquad
U''=\frac{6/3}{1}=\SI{2}{\volt};
\qquad
U=\SI{8}{\volt}.
\]
\textbf{Erro comum:} remover a fonte de $\SI{6}{\volt}$ \emph{junto} com seus
$\SI{3}{\ohm}$. Isso daria $G=\frac12+\frac16=\SI{0.667}{\siemens}$ e
$U'=\SI{9}{\volt}$ --- resultado inconsistente, pois a soma não fecharia com a
análise nodal.}
\end{questao}

\begin{questao}[2]{}
Verifique o resultado da questão anterior por análise nodal e explique por que
os dois métodos devem concordar.
\resposta{$U=\SI{8}{\volt}$ pelos dois métodos}
\resolucao{\textbf{Nodal:}
\[
\frac{U-12}{2}+\frac{U-6}{3}+\frac{U}{6}=0.
\]
Multiplicando por $6$: $3U-36+2U-12+U=0\Rightarrow6U=48\Rightarrow
U=\SI{8}{\volt}$ \checkmark.\\
\textbf{Por que concordam:} a equação nodal tem a forma
$G\,U=\sum b_k$, com $b_k$ as contribuições de corrente de cada fonte. Como
$U=\left(\sum b_k\right)/G=\sum\left(b_k/G\right)$, cada parcela $b_k/G$ é
exatamente a contribuição individual da fonte $k$.\\
A superposição não é um método \emph{alternativo} à análise nodal --- ela é a
mesma equação, lida termo a termo. A demonstração formal está no bloco de
desafio.}
\end{questao}

\blocoaprofundamento

\begin{questao}[3]{}
Um circuito contém uma fonte independente de $\SI{12}{\volt}$ (através de
$\SI{2}{\kilo\ohm}$ até o nó $A$), $\SI{2}{\kilo\ohm}$ de $A$ ao terra, uma
fonte de corrente independente de $\SI{1}{\milli\ampere}$ injetando em $A$, e
uma fonte \textbf{dependente} que injeta $gV_A$ em $A$, com
$g=\SI{0.5}{\milli\siemens}$.
\begin{itensq}
\item Determine $V_A$ pela análise nodal.
\item Determine-o por superposição, aplicando o procedimento correto.
\item Mostre o resultado que se obteria desativando a fonte dependente.
\end{itensq}
\resposta{(a) e (b) $V_A=\SI{14}{\volt}$; (c) $\SI{7}{\volt}$ --- errado}
\resolucao{(a) LKC em $A$:
\[
\frac{V_A-12}{2000}+\frac{V_A}{2000}-\pot{0.5}{-3}V_A=\pot{1}{-3}.
\]
A condutância líquida é
$\pot{1}{-3}-\pot{0.5}{-3}=\pot{0.5}{-3}\,\si{\siemens}$, e o lado direito vale
$\pot{6}{-3}+\pot{1}{-3}=\pot{7}{-3}$. Logo
$V_A=\SI{14}{\volt}$.\\
(b) \textbf{A fonte dependente permanece ativa nas duas análises parciais} ---
ela não é excitação, é parte da rede. A condutância líquida é sempre
$\pot{0.5}{-3}\,\si{\siemens}$:
\[
V'=\frac{\pot{6}{-3}}{\pot{0.5}{-3}}=\SI{12}{\volt};
\qquad
V''=\frac{\pot{1}{-3}}{\pot{0.5}{-3}}=\SI{2}{\volt};
\qquad
V_A=\SI{14}{\volt}.\ \checkmark
\]
(c) Desativando indevidamente a dependente, a condutância voltaria a
$\pot{1}{-3}$:
\[
V'=6,\quad V''=1
\;\Longrightarrow\;
V_A=\SI{7}{\volt},
\]
exatamente \textbf{metade} do correto.\\
\textbf{Critério:} desative apenas fontes cujo valor é \emph{imposto de fora}.
Se o valor depende de uma variável do próprio circuito, a fonte fica.}
\end{questao}

\begin{questao}[3]{}
Escreva $V_A$ do circuito de três fontes na forma
$V_A=a\,V_1+b\,V_2+c\,I_3$ e determine os coeficientes.
\begin{itensq}
\item Obtenha $a$, $b$ e $c$ com suas unidades.
\item Verifique com os valores originais.
\item Explique o significado físico de cada coeficiente.
\end{itensq}
\resposta{$a=1/6$, $b=1/3$ (adimensionais), $c=\SI{2}{\ohm}$}
\resolucao{(a) Com $G=\SI{0.5}{\siemens}$, cada fonte de tensão $V_k$ contribui
com $(V_k/R_k)/G$ e a fonte de corrente com $I_3/G$:
\[
a=\frac{1/12}{0{,}5}=\frac16;
\qquad
b=\frac{1/6}{0{,}5}=\frac13;
\qquad
c=\frac{1}{0{,}5}=\SI{2}{\ohm}.
\]
(b) $\frac{24}{6}+\frac{12}{3}+2(1)=4+4+2=\SI{10}{\volt}$ \checkmark.\\
(c) $a$ e $b$ são \textbf{ganhos de tensão} (adimensionais): que fração de cada
fonte chega ao nó. $c$ é uma \textbf{resistência de transferência}
($\si{\volt\per\ampere}$): quantos volts aparecem por ampère injetado.\\
\textbf{Utilidade:} obtidos uma única vez, os coeficientes respondem a qualquer
combinação de valores de fonte sem refazer o circuito. É a forma compacta de
armazenar a solução de uma rede linear --- e é exatamente a linha correspondente
da matriz $\mathbf{G}^{-1}$.}
\end{questao}

\begin{questao}[3]{}
Compare o esforço de resolver por superposição e por análise nodal um circuito
com $2$ nós desconhecidos e $4$ fontes independentes.
\begin{itensq}
\item Quantifique cada abordagem.
\item Indique quando cada uma é preferível.
\end{itensq}
\resposta{Superposição: 4 sistemas $2\times2$; nodal: 1 sistema $2\times2$}
\resolucao{(a) \textbf{Nodal:} monta-se \emph{um} sistema $2\times2$ com todas
as fontes no vetor independente e resolve-se uma vez.\\
\textbf{Superposição:} quatro sistemas $2\times2$, um por fonte, e depois a
soma. Cerca de \textbf{quatro vezes} mais trabalho.\\
(b) \textbf{Nodal é preferível} quando se quer apenas o resultado numérico ---
quase sempre.\\
\textbf{Superposição é preferível} quando: \emph{(i)} se quer saber a
\emph{contribuição} de cada fonte (análise de sensibilidade, rastreamento de
ruído, identificação da fonte dominante); \emph{(ii)} as fontes são de naturezas
diferentes e se deseja tratá-las separadamente; \emph{(iii)} algumas análises
parciais são triviais, colapsando o circuito por curtos ou aberturas ---
situação frequente e que pode inverter completamente a comparação de esforço.\\
\textbf{Síntese:} a superposição raramente é o caminho mais curto para o
número, mas frequentemente é o único caminho para o \emph{entendimento}.}
\end{questao}

\begin{questao}[3]{}
Um circuito tem duas fontes cujas contribuições para a corrente em um resistor
de $\SI{5}{\ohm}$ são $\SI{2}{\ampere}$ e $\SI{-3}{\ampere}$.
\begin{itensq}
\item Determine a corrente e a potência reais.
\item Calcule o erro que se cometeria somando potências parciais.
\item Determine o termo cruzado e seu sinal.
\end{itensq}
\resposta{(a) $\SI{-1}{\ampere}$, $\SI{5}{\watt}$; (b) somaria
$\SI{65}{\watt}$; (c) termo cruzado $=\SI{-60}{\watt}$}
\resolucao{(a) $I=2-3=\SI{-1}{\ampere}$, logo
$P=5(1)^{2}=\SI{5}{\watt}$.\\
(b) $P'=5(4)=20$ e $P''=5(9)=\SI{45}{\watt}$; a soma daria
$\SI{65}{\watt}$ --- \textbf{treze vezes} o valor real.\\
(c)
\[
P=R(I'+I'')^{2}=R I'^{2}+R I''^{2}+2RI'I''
\;\Longrightarrow\;
P_{cruz}=2(5)(2)(-3)=\SI{-60}{\watt}.
\]
\textbf{Sinal negativo:} as contribuições têm sentidos opostos e se
\emph{cancelam} parcialmente. O termo cruzado subtrai.\\
\textbf{Caso extremo:} se as contribuições fossem $+2$ e $\SI{-2}{\ampere}$, a
corrente real seria nula e a potência também --- enquanto a soma de potências
parciais daria $\SI{40}{\watt}$ de dissipação inexistente. É o exemplo mais
claro de que potência não se superpõe.}
\end{questao}

\begin{questao}[3]{}
Duas fontes de tensão de $\SI{5}{\volt}$ e $\SI{-5}{\volt}$, cada uma através
de $\SI{1}{\kilo\ohm}$, alimentam um nó ligado ao terra por um \textbf{diodo}
(queda direta de $\SI{0.7}{\volt}$, bloqueio perfeito no reverso).
\begin{itensq}
\item Determine a tensão do nó com cada fonte agindo sozinha.
\item Some as contribuições.
\item Determine a tensão real e explique a discrepância.
\end{itensq}
\resposta{(a) $\SI{0.7}{\volt}$ e $\SI{-2.5}{\volt}$; (b) previsão
$\SI{-1.8}{\volt}$; (c) real $\SI{0}{\volt}$}
\resolucao{(a) \textbf{Só a fonte $+\SI{5}{\volt}$} (a outra substituída por
curto): o nó é alimentado por $\SI{5}{\volt}$ através de
$\SI{1}{\kilo\ohm}$ e ligado ao terra por outro $\SI{1}{\kilo\ohm}$ (o ramo
curto-circuitado). Sem o diodo, o divisor daria $\SI{2.5}{\volt}$ --- acima da
queda direta. O diodo conduz e grampeia o nó em
$V'=\SI{0.7}{\volt}$.\\
\textbf{Só a fonte $-\SI{5}{\volt}$:} pelo mesmo divisor, o nó tenderia a
$\SI{-2.5}{\volt}$. O diodo fica reversamente polarizado e não conduz, de modo
que $V''=\SI{-2.5}{\volt}$.\\
(b) Soma das parcelas: $0{,}7-2{,}5=\SI{-1.8}{\volt}$.\\
(c) \textbf{Situação real:} com as duas fontes ativas e simétricas, as correntes
pelos dois resistores iguais se cancelam e o nó assume $\SI{0}{\volt}$. Com
$\SI{0}{\volt}$ o diodo não conduz --- hipótese consistente. Logo
$V=\SI{0}{\volt}$.\\
A previsão da superposição ($\SI{-1.8}{\volt}$) está \textbf{errada}, e nem
sequer perto.\\
\textbf{Razão:} o diodo não é linear --- seu \emph{estado} (conduzindo ou
bloqueado) depende do circuito inteiro, e mudou entre as duas análises
parciais. Superpor respostas obtidas em \emph{configurações diferentes} não tem
significado. Este é o contraexemplo definitivo: a superposição não é uma
conveniência que se pode forçar, é uma consequência da linearidade.}
\end{questao}

\begin{questao}[3]{}
Explique por que, ao desativar uma fonte \textbf{real} de tensão, mantém-se sua
resistência interna, e quantifique o erro de removê-la no circuito de
$E_1=\SI{12}{\volt}/r_1=\SI{2}{\ohm}$, $E_2=\SI{6}{\volt}/r_2=\SI{3}{\ohm}$ e
$R_L=\SI{6}{\ohm}$.
\resposta{Contribuição correta de $E_1$: $\SI{6}{\volt}$; removendo $r_2$
obter-se-ia $\SI{9}{\volt}$ --- erro de $50\%$}
\resolucao{\textbf{Por quê:} a fonte real é um \emph{modelo} de dois elementos
--- uma f.e.m. ideal em série com $r$. Desativar a \emph{fonte} significa zerar
a f.e.m.; o resistor $r$ é um componente físico da rede e permanece, exatamente
como qualquer outro resistor.\\
\textbf{Correto:} com $E_2=0$ e $r_2$ presente,
$G=\frac12+\frac13+\frac16=\SI{1}{\siemens}$ e
\[
U'=\frac{12/2}{1}=\SI{6}{\volt}.
\]
\textbf{Incorreto:} removendo o ramo inteiro,
$G=\frac12+\frac16=\SI{0.667}{\siemens}$ e
$U'=\dfrac{6}{0{,}667}=\SI{9}{\volt}$ --- $50\%$ acima.\\
\textbf{Como o erro se denuncia:} a soma das parcelas incorretas não fecha com
a análise nodal ($\SI{8}{\volt}$). Conferir a soma total contra um método
independente é a salvaguarda padrão da superposição.}
\end{questao}

\begin{questao}[3]{}
Em um circuito com duas fontes, mediu-se: com ambas ativas,
$I=\SI{1.4}{\ampere}$; com apenas a primeira, $I'=\SI{2.0}{\ampere}$.
\begin{itensq}
\item Determine a contribuição da segunda fonte.
\item Determine a corrente se a segunda fonte for dobrada.
\item Determine o valor da segunda fonte que anularia $I$.
\end{itensq}
\resposta{(a) $\SI{-0.6}{\ampere}$; (b) $\SI{0.8}{\ampere}$; (c) o triplo do
valor original}
\resolucao{(a) $I''=I-I'=1{,}4-2{,}0=\SI{-0.6}{\ampere}$ --- ela se opõe à
primeira.\\
(b) Dobrando a segunda fonte, sua contribuição dobra e a da primeira não muda:
\[
I=2{,}0+2(-0{,}6)=\SI{0.8}{\ampere}.
\]
(c) Anular exige $I''=\SI{-2.0}{\ampere}$, ou seja, um fator
$2{,}0/0{,}6=3{,}33$ sobre o valor original.\\
\textbf{Poder do método:} duas medições permitem prever o comportamento para
\emph{qualquer} valor da segunda fonte, sem conhecer a topologia do circuito.
Basta a garantia de linearidade --- e é justamente por isso que a verificação
experimental da linearidade (bloco de desafio) tem valor prático.}
\end{questao}

\begin{questao}[3]{}
Uma ponte é alimentada por duas fontes independentes. Sabe-se que a
contribuição da primeira para a corrente do galvanômetro é
$\SI{+40}{\milli\ampere}$ e que a corrente total medida é
$\SI{0}{\milli\ampere}$.
\begin{itensq}
\item Determine a contribuição da segunda fonte.
\item Explique por que o equilíbrio observado \emph{não} significa que a ponte
      esteja balanceada no sentido usual.
\end{itensq}
\resposta{(a) $\SI{-40}{\milli\ampere}$; (b) trata-se de cancelamento entre
fontes, não de equilíbrio da rede}
\resolucao{(a) $I''=0-40=\SI{-40}{\milli\ampere}$.\\
(b) O \textbf{equilíbrio clássico} da ponte é uma propriedade da \emph{rede}
($R_1/R_2=R_3/R_4$): a corrente é nula qualquer que seja a alimentação, e o
valor de $R_g$ é irrelevante.\\
Aqui, ao contrário, a rede está desequilibrada --- cada fonte, sozinha, produz
corrente considerável. O zero resulta de um \textbf{cancelamento numérico}
entre as duas contribuições, e depende dos \emph{valores} das fontes.\\
\textbf{Consequência prática:} basta variar qualquer uma das fontes em $1\%$
para que o zero desapareça. Um ajuste desse tipo não tem a robustez do
equilíbrio verdadeiro, e por isso não serve como método de medição --- ele
mediria a razão entre as fontes, não a razão entre os resistores.}
\end{questao}

\begin{questao}[3]{}
Um sinal de $\SI{50}{\milli\volt}$ e uma interferência de rede de
$\SI{200}{\milli\volt}$ atuam simultaneamente sobre um circuito linear, com
ganhos de $8$ e $0{,}05$ respectivamente.
\begin{itensq}
\item Determine as contribuições na saída.
\item Determine a razão sinal-ruído resultante.
\item Explique como a superposição justifica o tratamento separado.
\end{itensq}
\resposta{(a) $\SI{400}{\milli\volt}$ e $\SI{10}{\milli\volt}$; (b) $40$, ou
$\SI{32}{\decibel}$}
\resolucao{(a) $V_{sinal}=8(50)=\SI{400}{\milli\volt}$;
$V_{ruído}=0{,}05(200)=\SI{10}{\milli\volt}$.\\
(b) Razão $=400/10=40$, ou
$20\log_{10}40\approx\SI{32}{\decibel}$.\\
Na entrada, a razão era $50/200=0{,}25$ --- o \emph{ruído era quatro vezes
maior que o sinal}. O circuito melhorou a relação em um fator $160$.\\
(c) \textbf{Justificativa:} em rede linear, sinal e interferência percorrem o
circuito \emph{independentemente}, cada um com seu próprio ganho. Isso permite
projetar o circuito para maximizar um ganho e minimizar o outro, tratando-os
como problemas separados.\\
\textbf{E é exatamente isso que se perde na não linearidade:} em um circuito não
linear surgem \emph{produtos de intermodulação} --- termos cruzados que geram
componentes em frequências novas, impossíveis de separar depois. É o análogo,
em sinais, do termo cruzado de potência.}
\end{questao}

\begin{questao}[3]{}
Um circuito linear tem três fontes. Sabe-se que $V_{saída}=\SI{15}{\volt}$ com
todas ativas, $\SI{9}{\volt}$ sem a terceira e $\SI{4}{\volt}$ apenas com a
primeira. Determine a contribuição de cada fonte.
\resposta{$\SI{4}{\volt}$, $\SI{5}{\volt}$ e $\SI{6}{\volt}$}
\resolucao{Sejam $a$, $b$ e $c$ as contribuições.
\begin{itemize}
\item Todas ativas: $a+b+c=15$.
\item Sem a terceira: $a+b=9$ $\Rightarrow$ $c=15-9=\SI{6}{\volt}$.
\item Apenas a primeira: $a=\SI{4}{\volt}$ $\Rightarrow$
      $b=9-4=\SI{5}{\volt}$.
\end{itemize}
\textbf{Verificação:} $4+5+6=\SI{15}{\volt}$ \checkmark.\\
\textbf{Observação de método:} bastaram \emph{três} medições para decompor
completamente um sistema de três fontes, e nenhuma delas exigiu isolar a segunda
ou a terceira sozinhas. Em geral, $m$ medições \textbf{independentes} bastam
para $m$ fontes --- resultado que decorre de o sistema ser linear e as
contribuições, aditivas.}
\end{questao}

\blocodesafio

\begin{questao}[4]{}
\textbf{(Demonstração formal.)} Seja uma rede resistiva descrita por
$\mathbf{G}\mathbf{v}=\mathbf{b}$, com
$\mathbf{b}=\mathbf{b}_1+\mathbf{b}_2$ as contribuições de duas fontes.
\begin{itensq}
\item Prove a superposição a partir da linearidade do sistema.
\item Generalize para $m$ fontes.
\item Identifique onde a demonstração falha se $\mathbf{G}$ depender de
      $\mathbf{v}$.
\end{itensq}
\resposta{$\mathbf{v}=\mathbf{G}^{-1}\mathbf{b}_1+\mathbf{G}^{-1}\mathbf{b}_2$}
\resolucao{(a) Como $\mathbf{G}$ é não singular em rede passiva conectada,
\[
\mathbf{v}=\mathbf{G}^{-1}\mathbf{b}
=\mathbf{G}^{-1}(\mathbf{b}_1+\mathbf{b}_2)
=\underbrace{\mathbf{G}^{-1}\mathbf{b}_1}_{\mathbf{v}_1}
+\underbrace{\mathbf{G}^{-1}\mathbf{b}_2}_{\mathbf{v}_2}.
\]
Ora, $\mathbf{v}_1$ é exatamente a solução com \emph{apenas} a primeira fonte
ativa (as demais zeradas anulam suas parcelas em $\mathbf{b}$). Logo
$\mathbf{v}=\mathbf{v}_1+\mathbf{v}_2$: o teorema.\\
(b) Por indução imediata,
$\mathbf{v}=\sum_{k=1}^{m}\mathbf{G}^{-1}\mathbf{b}_k$. A distributividade do
produto matricial sobre a soma é tudo o que se usa.\\
(c) \textbf{Se $\mathbf{G}=\mathbf{G}(\mathbf{v})$}, a matriz muda de uma
análise parcial para outra --- e $\mathbf{G}_1^{-1}\mathbf{b}_1
+\mathbf{G}_2^{-1}\mathbf{b}_2$ não é solução de nada. Foi precisamente isso que
ocorreu no exemplo do diodo: a matriz do circuito com o diodo conduzindo é
diferente da matriz com ele bloqueado.\\
\textbf{Três leituras do resultado.} \emph{(i)} A superposição é uma
propriedade da \emph{álgebra linear}, não da eletricidade --- vale para qualquer
sistema linear. \emph{(ii)} Fontes dependentes entram em $\mathbf{G}$, não em
$\mathbf{b}$ --- por isso não são desativadas. \emph{(iii)} A potência é
quadrática em $\mathbf{v}$, e nenhuma função quadrática de uma soma é a soma das
funções.}
\end{questao}

\begin{questao}[4]{}
\textbf{(O termo cruzado.)} Para um resistor $R$ sujeito a contribuições de
corrente $I_1$ e $I_2$:
\begin{itensq}
\item Expanda $P$ e identifique o termo cruzado.
\item Determine em que condições ele é nulo, positivo ou negativo.
\item Mostre que a soma de potências parciais pode ser maior \emph{ou} menor
      que a potência real, com exemplos numéricos.
\end{itensq}
\resposta{$P=RI_1^{2}+RI_2^{2}+2RI_1I_2$; o cruzado tem o sinal de
$I_1I_2$}
\resolucao{(a)
\[
P=R(I_1+I_2)^{2}=\underbrace{RI_1^{2}}_{P_1}+\underbrace{RI_2^{2}}_{P_2}
+\underbrace{2RI_1I_2}_{P_{cruz}}.
\]
(b) $P_{cruz}=0$ apenas se $I_1=0$ ou $I_2=0$ --- ou seja, praticamente nunca
com duas fontes atuantes. Ele é \textbf{positivo} quando as contribuições têm o
mesmo sentido e \textbf{negativo} quando se opõem.\\
\emph{(Nota: em corrente alternada, $P_{cruz}$ pode ter média nula se as
contribuições forem de frequências diferentes --- é a única situação em que
somar potências é legítimo, e ela não ocorre em CC.)}\\
(c) \textbf{Soma parcial menor que a real} ($R=\SI{4}{\ohm}$, contribuições de
tensão $6$ e $\SI{2}{\volt}$):
$P_1+P_2=9+1=\SI{10}{\watt}$ contra $P=\SI{16}{\watt}$. O cruzado soma
$\SI{6}{\watt}$.\\
\textbf{Soma parcial maior que a real} ($R=\SI{5}{\ohm}$, correntes
$+2$ e $\SI{-3}{\ampere}$): $P_1+P_2=20+45=\SI{65}{\watt}$ contra
$P=\SI{5}{\watt}$. O cruzado subtrai $\SI{60}{\watt}$.\\
\textbf{Conclusão:} o erro não tem sequer sinal previsível --- não existe
"correção" ou fator de segurança que salve o procedimento. A única saída é
superpor tensões ou correntes e calcular a potência ao final.}
\end{questao}

\begin{questao}[4]{}
\textbf{(Algoritmo.)} Formule um procedimento de superposição para $m$ fontes e
$n$ grandezas de interesse.
\begin{itensq}
\item Descreva o algoritmo e a estrutura de dados resultante.
\item Determine o custo computacional e compare com $m$ análises completas.
\item Indique como a matriz resultante é usada depois.
\end{itensq}
\resposta{Matriz de coeficientes $n\times m$; custo de uma fatoração mais $m$
substituições}
\resolucao{(a) \textbf{Algoritmo.}
\begin{enumerate}
\item Monte $\mathbf{G}$ uma única vez (ela não depende das fontes).
\item Fatore $\mathbf{G}=\mathbf{L}\mathbf{U}$ --- também uma única vez.
\item Para cada fonte $k=1,\dots,m$: monte $\mathbf{b}_k$ com apenas essa fonte
      ativa e resolva por substituição, obtendo $\mathbf{v}_k$.
\item Extraia as $n$ grandezas de interesse de cada $\mathbf{v}_k$ e armazene
      na coluna $k$ de uma matriz $\mathbf{H}$ de dimensão $n\times m$.
\end{enumerate}
(b) \textbf{Custo:} a fatoração custa $\mathcal{O}(N^{3}/3)$ para uma rede de
$N$ nós, e cada substituição custa $\mathcal{O}(N^{2})$. Total:
$\mathcal{O}(N^{3}/3+mN^{2})$.\\
Repetir a análise completa $m$ vezes custaria
$\mathcal{O}(mN^{3}/3)$. Para $N=20$ e $m=5$, a razão é de aproximadamente
$4{,}7$ contra $13{,}3$ (em milhares de operações) --- \textbf{quase três vezes
mais rápido}, e a vantagem cresce com $m$.\\
\textbf{Observação:} a superposição sai \emph{de graça} como subproduto da
fatoração. O custo extra em relação a uma única análise é apenas
$(m-1)N^{2}$.\\
(c) \textbf{Uso posterior:} qualquer combinação de valores de fonte
$\mathbf{s}$ produz as saídas por um simples produto
$\mathbf{y}=\mathbf{H}\mathbf{s}$ --- sem tocar no circuito. A matriz
$\mathbf{H}$ é o \emph{modelo linear} completo da rede, e é ela que permite
varredura de parâmetros, análise de pior caso e otimização.}
\end{questao}

\begin{questao}[4]{}
\textbf{(Projeto de ensaio.)} Elabore uma verificação experimental da
superposição em um circuito com duas fontes de tensão.
\begin{itensq}
\item Descreva o procedimento e o cuidado essencial na desativação.
\item Estabeleça o critério de aceitação, dadas incertezas de $\pm1\%$ nas
      medições.
\item Interprete uma discrepância de $8\%$.
\end{itensq}
\resposta{Critério: $|V-(V'+V'')|$ dentro de cerca de $\pm3\%$; $8\%$ indica
não linearidade ou erro de procedimento}
\resolucao{(a) \textbf{Procedimento.}
\begin{enumerate}
\item Com ambas as fontes ativas, meça $V$ no ponto escolhido.
\item Desative a fonte $2$ e meça $V'$.
\item Restaure a fonte $2$, desative a $1$ e meça $V''$.
\item Compare $V$ com $V'+V''$.
\end{enumerate}
\textbf{Cuidado essencial:} desativar significa \textbf{substituir por um
curto}, não simplesmente desligar a fonte nem retirá-la do circuito. Uma fonte
desligada apresenta impedância indefinida (alta, em geral) --- o que equivale a
\emph{abrir} o ramo e altera a topologia. Se a fonte tiver resistência interna
apreciável, ela deve ser mantida no lugar da f.e.m.\\
(b) \textbf{Critério.} As três medições são independentes, e as incertezas se
propagam. Como $V'$ e $V''$ entram somadas, a incerteza da soma é
$\sqrt{u'^{2}+u''^{2}}$; combinada com a de $V$, chega-se a algo em torno de
$\pm1{,}7\%$ para uma estimativa estatística, ou até $\pm3\%$ no pior caso.
\textbf{Adote $\pm3\%$} como faixa de aceitação.\\
(c) \textbf{Discrepância de $8\%$} está muito além da incerteza. Causas
prováveis, em ordem de probabilidade:
\begin{itemize}
\item \textbf{Erro de procedimento:} fonte desligada em vez de curto-circuitada
      --- de longe a causa mais comum.
\item \textbf{Resistência interna ignorada} na desativação.
\item \textbf{Não linearidade real:} aquecimento de resistores, diodo de
      proteção conduzindo, ou fonte entrando em limitação.
\end{itemize}
\textbf{Como distinguir:} repita o ensaio com tensões \emph{reduzidas à
metade}. Se a discrepância percentual permanecer, o problema é de procedimento;
se diminuir, é não linearidade dependente de nível.}
\end{questao}

\begin{questao}[4]{}
\textbf{(Comparação de métodos.)} Um circuito tem duas fontes de tensão reais e
uma fonte de corrente.
\begin{itensq}
\item Descreva a solução por transformação de fontes.
\item Compare com a superposição em número de etapas e em informação obtida.
\item Indique o critério de escolha.
\end{itensq}
\resposta{Transformação: mais rápida para o número; superposição: entrega a
decomposição}
\resolucao{(a) \textbf{Transformação de fontes.} Converta cada fonte real de
tensão $(E_k,r_k)$ em fonte de corrente $(E_k/r_k, r_k)$. Todas as fontes de
corrente ficam então em paralelo com todas as resistências, e podem ser
\emph{somadas algebricamente} em uma única fonte equivalente; as condutâncias
também se somam. Uma divisão final entrega a tensão do nó.\\
(b) \textbf{Etapas.}
\begin{center}
\begin{tabular}{lcc}
\hline
& Transformação & Superposição\\
\hline
Análises do circuito & 1 & 3\\
Resultado numérico & sim & sim\\
Contribuição por fonte & não & sim\\
\hline
\end{tabular}
\end{center}
A transformação é claramente mais econômica --- ela é, na prática, a
\emph{montagem} da equação nodal por outro caminho.\\
Mas note: as parcelas $E_k/r_k$ que a transformação soma são \emph{exatamente}
as contribuições que a superposição calcula separadamente. Os métodos não são
rivais; um é a forma agregada do outro.\\
(c) \textbf{Critério.} Use transformação de fontes (ou análise nodal) para
obter o número. Use superposição quando a pergunta for \emph{"quanto cada fonte
contribui"} --- rastreamento de ruído, análise de sensibilidade, diagnóstico de
qual fonte está causando um desvio. Se a resposta desejada for um número único,
a superposição é trabalho desperdiçado.}
\end{questao}

\begin{questao}[4]{}
\textbf{(Limites de aplicação.)} Determine se a superposição se aplica e
justifique.
\begin{itensq}
\item Circuito com lâmpada incandescente, cuja resistência cresce com a
      temperatura.
\item Circuito linear em que se deseja a \emph{energia} dissipada em um
      intervalo.
\item Circuito com amplificador operacional operando na região linear.
\end{itensq}
\resposta{(a) não; (b) não; (c) sim}
\resolucao{(a) \textbf{Não se aplica.} A resistência depende da potência
dissipada, que depende da corrente --- realimentação que torna o circuito não
linear. Cada análise parcial teria uma temperatura, e portanto uma
\emph{resistência}, diferente.\\
\emph{Ressalva:} para pequenas variações em torno de um ponto de operação, o
modelo incremental é linear e a superposição vale \emph{para as variações}.\\
(b) \textbf{Não se aplica.} Energia é a integral da potência, e potência é
quadrática --- herda o termo cruzado. Superpor energias comete o mesmo erro que
superpor potências, apenas acumulado no tempo.\\
\emph{Procedimento correto:} superponha correntes, obtenha $i(t)$ total,
calcule $p(t)=Ri^{2}$ e só então integre.\\
(c) \textbf{Aplica-se.} O amp-op na região linear é descrito por relações
lineares (curto virtual e corrente de entrada nula), e o ganho é uma constante.
Superposição é, aliás, a ferramenta padrão para analisar amplificadores
somadores e diferenciais: calcula-se o ganho de cada entrada com as demais
aterradas e somam-se as contribuições.\\
\emph{Fronteira:} se a saída satura, a linearidade se perde e a superposição
falha --- o mesmo padrão do diodo.}
\end{questao}

\begin{questao}[4]{}
\textbf{(Escolha de sinais.)} Estabeleça um procedimento sistemático para os
sinais das contribuições e aplique-o a um caso com resultado negativo.
\begin{itensq}
\item Formule a regra.
\item Aplique a duas fontes cujas contribuições de corrente são
      $\SI{2}{\ampere}$ (a favor) e $\SI{5}{\ampere}$ (contra) a referência.
\item Explique por que o resultado negativo não indica erro.
\end{itensq}
\resposta{$I=\SI{-3}{\ampere}$: mesma intensidade, sentido oposto ao adotado}
\resolucao{(a) \textbf{Regra em três passos.}
\begin{enumerate}
\item \textbf{Antes} de qualquer cálculo, desenhe uma seta de referência para a
      corrente (ou uma polaridade $+/-$ para a tensão) e \textbf{não a mude
      mais}.
\item Em cada análise parcial, determine o sentido real da contribuição.
Atribua sinal $+$ se coincidir com a referência e $-$ se for oposto.
\item Some algebricamente.
\end{enumerate}
O erro clássico é redesenhar a seta a cada análise parcial "para ficar
positivo" --- o que torna a soma final sem significado.\\
(b)
\[
I=(+2)+(-5)=\SI{-3}{\ampere}.
\]
(c) \textbf{O sinal negativo é informação, não erro.} Ele diz que a corrente
real tem módulo $\SI{3}{\ampere}$ e flui no sentido \emph{oposto} ao da seta
desenhada. Se a seta tivesse sido desenhada ao contrário desde o início, o
resultado seria $\SI{+3}{\ampere}$ --- e descreveria exatamente a mesma
realidade física.\\
\textbf{Verificação útil:} grandezas que não podem ser negativas --- como a
potência dissipada em um resistor --- servem de teste. Aqui,
$P=R(3)^{2}$ é positiva independentemente do sinal de $I$, como deve ser.}
\end{questao}

\begin{questao}[4]{}
\textbf{(Superposição e ruído.)} Um amplificador linear recebe um sinal
$V_s$ e sofre a injeção de uma interferência $V_n$ em um ponto interno, com
ganhos $A_s=40$ e $A_n=2$.
\begin{itensq}
\item Escreva a saída e a razão sinal-ruído em função das entradas.
\item Determine o ganho $A_n$ máximo admissível para
      $\mathrm{SNR}\ge\SI{40}{\decibel}$, com
      $V_s=\SI{10}{\milli\volt}$ e $V_n=\SI{100}{\milli\volt}$.
\item Explique por que essa análise seria impossível em circuito não linear.
\end{itensq}
\resposta{(a) $V_o=A_sV_s+A_nV_n$; (b) $A_n\le0{,}04$}
\resolucao{(a) Pela superposição, cada entrada percorre o circuito
independentemente:
\[
V_o=A_sV_s+A_nV_n;
\qquad
\mathrm{SNR}=\frac{A_sV_s}{A_nV_n}.
\]
(b) $\SI{40}{\decibel}$ corresponde a uma razão de $100$:
\[
\frac{40(10)}{A_n(100)}\ge100
\;\Longrightarrow\;
\frac{400}{100A_n}\ge100
\;\Longrightarrow\;
A_n\le0{,}04.
\]
O ganho de $2$ do enunciado é \textbf{cinquenta vezes} maior que o admissível
--- o projeto está muito longe da especificação, e a conta diz exatamente
quanto: são necessários $\SI{34}{\decibel}$ de atenuação adicional no caminho
da interferência.\\
(c) \textbf{Em circuito não linear, "o ganho do ruído" não existe como grandeza
independente.} A resposta à interferência dependeria do nível do sinal
\emph{simultâneo}, e vice-versa. Surgiriam produtos de intermodulação --- termos
proporcionais a $V_sV_n$ --- que criam componentes espectrais novas, dentro da
banda do sinal e portanto \textbf{impossíveis de filtrar}.\\
É o mesmo termo cruzado que impede somar potências, agora manifestado no
domínio dos sinais. A linearidade não é um detalhe matemático: é o que garante
que sinal e ruído permaneçam separáveis.}
\end{questao}
